\section{Related Work}
In this section, we review the related work for ITE estimation using static and time-varying observational data. We first introduce the causal effect learning framework on static data, and then the framework based on time-varying data. 

\noindent\textbf{Learning causal effects with static data} According to the way to control the confounders, existing work with static observational data can be divided into four groups: 1) Matching-based methods; 2) Tree-based methods; 3) Reweighting-based methods; 4) Representation-based methods. The matching-based methods are adopted to estimate the counterfactual from the nearest neighbors. The distances among individuals can be measured in several ways (i.e., Euclidean distance, propensity scores). For example, propensity score matching (PSM) \cite{rosenbaum1983central} is to match a treated (control) sample to a set of control (treated) samples with similar propensity scores. Tree-based methods are also widely adopted in causal effect estimation. Bayesian additive regression trees (BART) \cite{hill2011bayesian} is a non-parametric Bayesian regression tree model based on the \textit{strong ignorability assumption}. It is easy to implement and free from parameter adjustment. Causal Forest (CF) \cite{wager2018estimation} is also a tree-based causal effect estimation method, which estimates the treatment effect at the leaf node by mapping the original covariate into tree and forests. The reweighting-based methods attempt to re-weight samples in the population for correcting the bias in observational data. For example, inverse probability of treatment weighting (IPTW) \cite{rosenbaum1983central} removes the confounding by assigning a weight to each individual in the population. The weights are calculated based on the propensity score. Recently, representation learning methods are proposed for causal effect estimation via balancing the distribution between treated and control groups in hidden space \cite{johansson2016learning,shalit2017estimating}. Moreover, Yao et al. \cite{yao2018representation} incorporate the local similarity among individuals with population-level distribution balancing in latent space to better estimate ITE. Yoon et al. propose to use generative adversarial nets (GAN) for inferring the counterfactual outcomes based on factual outcomes. Shi et al. \cite{shi2019adapting} jointly model the propensity prediction and potential outcome prediction as a multi-task learning problem. 

Though existing work shows great performance in causal effect estimation, they still have some limitations. First, most of them are built upon \textit{strong ignorability assumption} without considering the influence of hidden confounders. This constrain has been shown to lead to bias in estimating causal effects \cite{pearl2009causality}. Moreover, existing work is initially designed for static data, which is not easy to adapt for ITE estimation under dynamic longitudinal setting. 

\vspace{3pt}
\noindent\textbf{Learning causal effects with time-varying data}
As estimating causal effect from observational data is significant and most observational data contains sequential information, some work has been proposed for dealing time-varying confounders. In statistics and epidemiology domains, a group of methods use the inverse probability of treatment and \textit{g-formula} based method to estimate causal effect with sequential data \cite{robins2000marginal,schulam2017reliable}. More recently, Lim et al. \cite{lim2018forecasting} propose a recurrent marginal structural network for predicting the patient's potential response to a series of treatments. Bica et al. \cite{bica2020estimating} adopt adversarial training techniques to balance the historical confounding variables. Their method is based on the \textit{strong ignorability assumption}. Later, Bica et al. \cite{bica2019time} relax the assumption on strong ignorability and propose to estimate the treatment response with the existence of hidden confounders. 